%-------------------------
% EECS Resume in Latex
% Author : Eric Zhong
% Based off of : Jake's Resume - https://www.overleaf.com/latex/templates/jakes-resume/syzfjbzwjncs
% License : MIT
%------------------------

% Preserve this resume's layout across newer LaTeX installations.
\RequirePackage[2024/11/01]{latexrelease}
\documentclass[letterpaper,10pt]{article}

\usepackage[T1]{fontenc}
\usepackage{latexsym}
\usepackage[empty]{fullpage}
\usepackage{titlesec}
\usepackage{marvosym}
\usepackage[usenames,dvipsnames]{color}
\usepackage{verbatim}
\usepackage{enumitem}
\usepackage[hidelinks]{hyperref}
\usepackage{fancyhdr}
\usepackage[english]{babel}
\usepackage{tabularx}
\usepackage{fontawesome5}
\usepackage{multicol}
\usepackage{siunitx}
\setlength{\multicolsep}{-3.0pt}
\setlength{\columnsep}{-1pt}
\input{glyphtounicode}


%----------FONT OPTIONS----------
% sans-serif
% \usepackage[sfdefault]{FiraSans}
% \usepackage[sfdefault]{roboto}
 \usepackage[sfdefault]{noto-sans} %Good Option
% \usepackage[default]{sourcesanspro} %Good Option

% serif
% \usepackage{CormorantGaramond}
% \usepackage{charter}


\pagestyle{fancy}
\fancyhf{} % clear all header and footer fields
\fancyfoot{}
\renewcommand{\headrulewidth}{0pt}
\renewcommand{\footrulewidth}{0pt}

% Adjust margins
\addtolength{\oddsidemargin}{-0.6in}
\addtolength{\evensidemargin}{-0.6in}
\addtolength{\textwidth}{1.19in}
\addtolength{\topmargin}{-.75in}
\addtolength{\textheight}{1.55in} % Bottom margin (effectively)

\urlstyle{same}

\raggedbottom
\raggedright
\setlength{\tabcolsep}{0in}

% Sections formatting
\titleformat{\section}{
  \vspace{1pt}\scshape\raggedright\Large\bfseries
}{}{0em}{}[\color{black}\titlerule \vspace{-5pt}]

% Ensure that generate pdf is machine readable/ATS parsable
\pdfgentounicode=1

%-------------------------
% Custom commands
\newcommand{\resumeItem}[1]{
  \item \small{#1}\vspace{-2pt}
}

\newcommand{\classesList}[4]{
    \item\small{
        {#1 #2 #3 #4 \vspace{-2pt}}
  }
}

\newcommand{\resumeSubheading}[4]{
  \vspace{-2pt}\item
    \begin{tabular*}{1.0\textwidth}[t]{l@{\extracolsep{\fill}}r}
      \textbf{#1} & \textbf{\small #2} \\
      \small #3 & \small #4 \\
    \end{tabular*}\vspace{-7pt}
}

% Company block: company name and location, followed by one or more roles.
\newcommand{\resumeCompanyHeading}[2]{
  \vspace{-2pt}\item
    \begin{tabular*}{\textwidth}[t]{l@{\extracolsep{\fill}}r}
      {\fontsize{12}{12}\selectfont\textbf{#1}} & \small #2 \\
    \end{tabular*}\par\nobreak\vspace{-\parskip}\vspace{-4pt}
}

% Role block: title and dates, followed by its experience bullet list.
\newcommand{\resumeRoleHeading}[2]{
    \par\nobreak\vspace{4pt}\noindent
    \begin{tabular*}{\textwidth}[t]{l@{\extracolsep{\fill}}r}
      \hspace{0.2cm}\textbf{\normalsize #1} & \small #2 \\
    \end{tabular*}\vspace{-7pt}
}

\newcommand{\resumeProjectHeading}[2]{
    \item
    \begin{tabular*}{1.001\textwidth}{l@{\extracolsep{\fill}}r}
      \small#1 & \textbf{\small #2}\\
    \end{tabular*}\vspace{-7pt}
}

\newcommand{\resumeSubItem}[1]{\resumeItem{#1}\vspace{-4pt}}

\renewcommand\labelitemi{$\vcenter{\hbox{\tiny$\bullet$}}$}
\renewcommand\labelitemii{$\vcenter{\hbox{\tiny$\bullet$}}$}

\newcommand{\resumeSubHeadingListStart}{\begin{itemize}[leftmargin=0.0in, label={}]}
\newcommand{\resumeSubHeadingListEnd}{\end{itemize}}
\newcommand{\resumeItemListStart}{\begin{itemize}}
\newcommand{\resumeExperienceItemListStart}{\begin{itemize}}
\newcommand{\resumeItemListEnd}{\end{itemize}\vspace{-5pt}}

%-------------------------------------------
%%%%%%  RESUME STARTS HERE  %%%%%%%%%%%%%%%%%%%%%%%%%%%%


\begin{document}

%----------HEADING----------
% \begin{tabular*}{\textwidth}{l@{\extracolsep{\fill}}r}
%   \textbf{\href{http://sourabhbajaj.com/}{\Large Sourabh Bajaj}} & Email : \href{mailto:sourabh@sourabhbajaj.com}{sourabh@sourabhbajaj.com}\\
%   \href{http://sourabhbajaj.com/}{http://www.sourabhbajaj.com} & Mobile : +1-123-456-7890 \\
% \end{tabular*}

\begin{center}
    \textbf{\Huge \scshape Eric Zhong} \\ \vspace{2pt}
    \small \raisebox{-0.0\height}\faPhone\ 267-481-6608 ~ \href{mailto:eric@erictzhong.com}{\raisebox{-0.1\height}\faEnvelope { eric@erictzhong.com}} ~
    \href{https://linkedin.com/in/erictzhong}{\raisebox{-0.03\height}\faLinkedin\ {linkedin.com/in/erictzhong }}
    \href{https://github.com/etzhong}{\raisebox{-0.1\height}\faGithub\ {github.com/etzhong }}
    \href{https://www.erictzhong.com}{\raisebox{-0.1\height}\faGlobe\ {erictzhong.com}}
\end{center}


%-----------EDUCATION-----------
\vspace{-21pt}
\section{Education}
  \resumeSubHeadingListStart

    % Degrees
    \vspace{1pt}
    \resumeSubheading
      {\large{University of Southern California}}{GPA: 4.00\vspace{2pt}}
      {\textbf{M.S. Electrical Engineering} (Computer Architecture and Chip Design Specialization)}{Graduated May 2025}
    \resumeSubheading
      {}{\vspace{-12pt}}
      {\textbf{B.S. Computer Engineering and Computer Science}}{(B.S. + M.S.)}

    % Awards
    \vspace{-9pt}
    \resumeItem{{\textbf{Honors}: Top of Class, Presidential Merit Scholarship, Provost Research Fellowship, Engineering Honors Fellowship}}

    % Coursework
    \vspace{-1pt}
    \resumeSubheading
        {{Relevant Coursework}}{}{
          \footnotesize{\hspace{0.58cm}\textbf{CS:} Algorithms and Computing Theory, Operating Systems, Compilers, Networking, C++, Data Structures and OOP, Software Engineering}\\
          \footnotesize{\hspace{0.5cm}\textbf{ECE:} ASIC Design and Tapeout, Computer Microarchitecture, Computer Systems Architecture, VLSI Circuit Design, System-on-Chip (SoCs),}\\
            \footnotesize{\hspace{1.07cm} FPGAs, Digital System Design, Network Processor Design, Parallel Computing Systems, Embedded and Distributed Systems}\\
          \footnotesize{\hspace{0.25cm}\textbf{Math:} Probability and Statistics, Linear Algebra, Discrete Math, Multivariable Calculus, Differential Equations}
        }{}
  \resumeSubHeadingListEnd


%-----------EXPERIENCE-----------
\vspace{-12pt}
\section{Experience}
  \resumeSubHeadingListStart

    \resumeCompanyHeading
      {NVIDIA}{\textbf{Santa Clara, CA}}
    \resumeRoleHeading
      {GPU Architect}{\textbf{June 2025---Present}}
    \resumeExperienceItemListStart
        \resumeItem{Driving hardware design, architecture specification, \textbf{C++} simulation, and functional verification for NVIDIA's next-generation Feynman GPUs. Owned IPs include host \textbf{DMA engines} and work schedulers, \textbf{PCIe}/\textbf{CXL} endpoints, and interrupt controllers}
        \resumeItem{Writing custom firmware in \textbf{C} targetting \textbf{NV-RISC-V assembly} for on-chip dynamic CUDA work scheduling and dependency resolution---enabling low-latency GPU kernel launch and improved performance with CUDA graph workflows}
        \resumeItem{Developing simulation tools in \textbf{C++} and \textbf{C} within NVIDIA's proprietary resman kernel driver and MODS testing framework for architectural studies and bring-up of new GPU hardware + software features}
        \resumeItem{Performing functional and performance testing using a \textbf{Python DSL} across \textbf{C++} simulation, \textbf{RTL}, \textbf{FPGA} emulation, and silicon}
        \resumeItem{Creating a new \textbf{Python} and \textbf{Perl} based dev flow for communicating multi-die hardware device information to driver software}
        \resumeItem{Training and tuning a new \textbf{AI agent} for \textbf{DFD/DFT}---enabling autonomous debug and RCA of functional failures on end silicon}
    \resumeItemListEnd
    \vspace{-1pt}
    \resumeRoleHeading
      {GPU Architecture Intern}{\textbf{May 2024---August 2024}}
      \resumeExperienceItemListStart
            \resumeItem{Architected and developed a \textbf{C++} functional model (fmodel) for a new GPU hardware IP enabling low-latency, real-time interrupt processing and work scheduling---facilitating RTL and kernel driver development for NVIDIA’s flagship Rubin GPUs}
            \resumeItem{Built and deployed a new lightweight unit testing infrastructure using \textbf{C++} and \textbf{gtest} within NVIDIA's proprietary simulation environment---accelerating functional and performance testing by up to a 20x speedup over existing fullchip testing workflows}
            \resumeItem{Refactored legacy \textbf{C++}/\textbf{SystemC} fmodel with \textbf{C++20} source to improve core sim performance and future code maintainability}
      \resumeItemListEnd
    \vspace{6pt}

    \resumeCompanyHeading
      {USC Photonic Computing Lab (Opticore)}{\textbf{Los Angeles, CA}}
    \resumeRoleHeading
      {Research Fellow}{\textbf{August 2022---May 2024}}
      \resumeExperienceItemListStart
        \resumeItem{Co-invented and published a novel teraFLOP optical tensor processor with a 25 fJ/Op energy efficiency---offering up to a 40x efficiency improvement over current state-of-the-art electronic processors for machine learning workflows (\href{https://www.science.org/doi/10.1126/sciadv.adu0228}{\underline{Publication Link}})}
        \resumeItem{Designed and verified a high-throughput (10+ GSPS), low-latency (< 100 ns) \textbf{FPGA}-based waveform generation and signal processing module using \textbf{SystemVerilog} and \textbf{C++}. Synthesized design onto an AMD Zynq UltraScale+ RFSoC using \textbf{Vivado}}
        \resumeItem{Developed a custom \textbf{Python} library for host interfacing with lab \textbf{FPGAs} along \textbf{100G Ethernet} to facilitate optical experiments}
      \resumeItemListEnd
    \vspace{6pt}

    \resumeCompanyHeading
      {IQVIA}{\textbf{Philadelphia, PA}}
    \resumeRoleHeading
      {Software Engineering Intern}{\textbf{May 2022---August 2022}}
      \resumeExperienceItemListStart
        \resumeItem{Engineered a high-throughput data pipeline using \textbf{Scala} and \textbf{Python} to process incoming market and research data feeds for 300+ healthcare products---aggregating and consolidating 40+ independent data sources for uniform ingestion and analysis}
        \resumeItem{Prototyped several multivariate regression models in \textbf{Python} using \textbf{Jupyter Notebook} for market forecasting across 20+ sectors---achieving an aggregate performance of <5\% SMAPE across a 3-month production testing window}
        \resumeItem{Built and deployed a \textbf{CI/CD} pipeline for \textbf{Scala microservices} backed by \textbf{Oracle DB}---automating pre-prod testing and SQA}
      \resumeItemListEnd

  \resumeSubHeadingListEnd

%-----------PROJECTS-----------
\vspace{-16pt}
\section{Projects}
    \vspace{-6pt}
    \resumeSubHeadingListStart

    % Architecture/Digital Design/RTL
    \resumeProjectHeading
    {\normalsize{\textbf{Cryptographic Accelerator ASIC}}}{}
    \resumeItemListStart
        \resumeItem{Developed architecture and \textbf{SystemVerilog RTL} for a power and area-optimized ASIC for efficiently computing variable delay functions (VDFs)---achieving 30+ million 256-bit hash computations per second with a module size of 105 um x 105 um}
        \resumeItem{Taped out design using Intel's proprietary 16nm FinFET process in February 2025. Design brought up in November 2025\\}
    \resumeItemListEnd

    % Architecture/Digital Design/RTL
    \vspace{-13pt}
    \resumeProjectHeading
    {\normalsize{\textbf{NetFPGA SmartNIC}}}{}
    \resumeItemListStart
        \resumeItem{Designed a high-throughput network switch for smart packet routing and preprocessing by payload content on a \textbf{NetFPGA-1G}}
        \resumeItem{Architected core IP and wrote \textbf{Verilog RTL} for a multi-core NPU, custom network routing logic, and a software-programmable memory interface---achieving 2.5M packet inferences per second. Verified and synthesized design using \textbf{Verilator} and \textbf{Vivado}}
    \resumeItemListEnd

    % Architecture/Digital Design/RTL
    \vspace{-13pt}
    \resumeProjectHeading
    {\normalsize{\textbf{Multithreaded ARM Processor}}}{}
    \resumeItemListStart
        \resumeItem{Designed a 32-bit, 4-way multithreaded, 8-stage pipelined processor in \textbf{Verilog} running a custom ARMv4-compatible ISA}
        \resumeItem{Synthesized and implemented \textbf{RTL} onto a Xilinx Virtex-II \textbf{FPGA} using \textbf{Vivado}. Verified design using \textbf{Verilator} and \textbf{cocotb}}
        \resumeItem{Built a custom Clang-based \textbf{compiler} toolchain and backend for converting \textbf{C} code into hardware-executable \textbf{assembly}}
    \resumeItemListEnd

    % Digital Design/RTL
    % \vspace{-12pt}
    % \resumeProjectHeading
    % {\normalsize{\textbf{FPGA Dinosaur Game}}}{}
    % \resumeItemListStart
    %     \resumeItem{Created a \textbf{Verilog} RTL version of the Google Chrome Dinosaur Game on an Artix-7 \textbf{FPGA} with animated VGA graphics}
    %     \resumeItem{Verified design using \textbf{Verilator} with a \textbf{C++} and \textbf{gtest} integrated functional model before synthesizing with \textbf{Vivado}}
    % \resumeItemListEnd

    \resumeSubHeadingListEnd

%-----------TECHNICAL SKILLS-----------
\vspace{-16pt}
\section{Skills}
 \begin{itemize}[leftmargin=0.0in, label={}]
    \vspace{-2pt}
    \small{\item{
     \textbf{Languages}{: C++, C, SystemVerilog, Verilog, Python, SystemC, Assembly, VHDL, Bash, Perl, Tcl} \\ %\hspace{1.75cm}%
     \textbf{Tools/Techs}{: FPGAs, Vivado, Vitis, Verilator, cocotb, ModelSim, Icarus Verilog, Quartus, Xilinx ISE, CUDA, LLVM, Linux, Git, P4 Helix} \\ %\hspace{0.48cm}
     \textbf{Domains}{: Hardware + Software Co-design, Computer Architecture, Systems Software, Open-Source Silicon, GPUs} \\%\hspace{0.48cm}
    }}
 \end{itemize}


\end{document}
https://www.overleaf.com/project/647970f4d4e9184c9e54ebed
